\section*{Bab \ref{ch:g8:balanced-equations} --- Kekekalan Massa dan Persamaan Reaksi Setara}

\begin{solution}{exo:g8:balanced-equations:1}
Dalam \omterm{def:g7:chemical-reactions:reaction}{reaksi kimia}, massa total produk yang terbentuk sama dengan massa
total \omterm{def:g7:chemical-reactions:reactant}{pereaksi} yang terpakai.
\end{solution}

\begin{solution}{exo:g8:balanced-equations:2}
Koefisiennya 3. Koefisien itu menyatakan 3 \omterm{def:g7:atoms-and-molecules:atom}{atom} karbon dan
$3 \times 2 = 6$ \omterm{def:g7:atoms-and-molecules:atom}{atom} oksigen.
\end{solution}

\begin{solution}{exo:g8:balanced-equations:3}
Ya: 1 C dan 2 O di setiap ruas.
\end{solution}

\begin{solution}{exo:g8:balanced-equations:4}
\ce{2Mg + O2 -> 2MgO}: 2 Mg dan 2 O di setiap ruas.
\end{solution}

\begin{solution}{exo:g8:balanced-equations:5}
Nitrogen: 2 N di kiri, jadi 2 \ce{NH3}; lalu ada 6 H di kanan, jadi 3
\ce{H2}: \ce{N2 + 3H2 -> 2NH3}.
\end{solution}

\begin{solution}{exo:g8:balanced-equations:6}
Menurut \omterm{def:g8:balanced-equations:conservation}{kekekalan massa}: $44 - 12 = \qty{32}{g}$ oksigen.
\end{solution}

\begin{solution}{exo:g8:balanced-equations:7}
Gas yang terbentuk lolos ke ruangan, sehingga neraca tidak lagi
menimbangnya. Tidak ada massa yang dimusnahkan: massa yang hilang itu
berada di \omterm{def:g4:air-a-mixture-of-gases:air}{udara} ruangan.
\end{solution}

\begin{solution}{exo:g8:balanced-equations:8}
Karbon: 1 di setiap ruas. Hidrogen: 4 di kiri, jadi 2 \ce{H2O}. Oksigen:
$2 + 2 = 4$ di kanan, jadi 2 \ce{O2}:
\ce{CH4 + 2O2 -> CO2 + 2H2O}.
\end{solution}

\begin{solution}{exo:g8:balanced-equations:9}
Mengubah \ce{H2O} menjadi \ce{H2O2} mengubah zatnya: \ce{H2O2} adalah
hidrogen peroksida, bukan air. Hanya koefisien yang boleh diubah:
\ce{2H2 + O2 -> 2H2O}.
\end{solution}

\begin{solution}{exo:g8:balanced-equations:10}
$20 \times 5 = 100$ \omterm{def:g7:atoms-and-molecules:molecule}{molekul} oksigen; $20 \times 4 = 80$ \omterm{def:g7:atoms-and-molecules:molecule}{molekul} air.
\end{solution}

\begin{solution}{exo:g8:balanced-equations:11}
Dengan 2 \ce{C2H6}: 4 C, jadi 4 \ce{CO2}; 12 H, jadi 6 \ce{H2O}; oksigen
di kanan $8 + 6 = 14$, jadi 7 \ce{O2}:
\ce{2C2H6 + 7O2 -> 4CO2 + 6H2O}.
\end{solution}

\begin{solution}{exo:g8:balanced-equations:12}
Tetap \qty{850.0}{g}. Kotak itu tertutup rapat: \qty{0.4}{g} lilin yang
\omterm{def:g4:what-a-fire-needs:burning}{terbakar} telah menjadi karbon dioksida dan uap air, yang masih berada di
dalam kotak bersama oksigen yang diambilnya. Tidak ada yang masuk atau
keluar.
\end{solution}

\begin{solution}{pb:g8:balanced-equations:1}
\textbf{1.} Dari oksigen \omterm{def:g4:air-a-mixture-of-gases:air}{udara}, yang bergabung dengan besi: oksida besi
seberat besi ditambah oksigen itu.

\textbf{2.} Semuanya tetap di dalam stoples: oksigen yang terpakai sudah
berada di dalam stoples, dan oksida yang terbentuk tetap di sana. Massa
totalnya tidak berubah.

\textbf{3.} \omterm{def:g8:balanced-equations:conservation}{Kekekalan massa}: dalam reaksi, massa total produk sama
dengan massa total \omterm{def:g7:chemical-reactions:reactant}{pereaksi}.

\textbf{4.} \ce{4Fe + 3O2 -> 2Fe2O3}.

\textbf{5.} \ce{3Fe + 2O2 -> Fe3O4}.

\textbf{6.} Kiri: 4 Fe, $3 \times 2 = 6$ O. Kanan: $2 \times 2 = 4$ Fe,
$2 \times 3 = 6$ O. Setara.

\textbf{7.} 4 \omterm{def:g7:atoms-and-molecules:atom}{atom} besi untuk 3 \omterm{def:g7:atoms-and-molecules:molecule}{molekul} oksigen, jadi $400 \times
\frac{3}{4} = 300$ \omterm{def:g7:atoms-and-molecules:molecule}{molekul} oksigen.

\textbf{8.} $0.28 \times \frac{3}{7} = \qty{0.12}{g}$ oksigen.

\textbf{9.} $0.28 + 0.12 = \qty{0.40}{g}$ oksida besi.

\textbf{10.} Ya: yang diperlukan \qty{0.12}{g} dan stoples itu berisi
sekitar \qty{0.5}{g}.

\textbf{11.} Oksigen yang terpakai telah meninggalkan \omterm{def:g4:air-a-mixture-of-gases:air}{udara} di dalam
stoples (sekarang berada di dalam oksida padat), sehingga stoples berisi
lebih sedikit gas daripada sebelumnya dan \omterm{def:g4:air-a-mixture-of-gases:air}{udara} dari ruangan mengalir
masuk. Angka neraca lalu naik, sebesar massa \omterm{def:g4:air-a-mixture-of-gases:air}{udara} yang masuk.

\textbf{12.} \qty{0.12}{g} oksigen diambil dari \omterm{def:g4:air-a-mixture-of-gases:air}{udara} di dalam stoples.
\end{solution}
